\documentclass{llncs}
\usepackage[utf8]{inputenc}
\usepackage[T1]{fontenc}
\usepackage{csquotes}
\usepackage{hyphenat}
\hyphenation{pro-ble-ma}

\usepackage{amssymb}
\usepackage{amsmath}
\usepackage{array}
\usepackage{graphicx}
\graphicspath{ {images/} }
\usepackage{listings}
\usepackage{booktabs}
\usepackage{multirow}

\lstset{
  basicstyle=\small\ttfamily,
  breaklines=true,
  frame=single,
  language=C
}

% Note: llncs.cls already provides theorem/lemma/corollary/proposition and a
% proof environment with \qed.  Do not \newtheorem them again.

\begin{document}

\title{Decentralized Federated Learning\\with zk-SNARKs and Byzantine-Resilient Consensus}

\author{Abderraouf LAARBI \and Kenza RIAHI}
\institute{ESILV\\
\email{abderraouf.laarbi@edu.devinci.fr, kenza.riahi@devinci.fr}}

\maketitle

% ─────────────────────────────────────────────────────────────────────────────
\begin{abstract}
Federated Learning lets participants train a shared model without pooling their
data, but in a decentralized setting nothing forces a participant to have
actually trained. We present a peer-to-peer federated learning system in which
each node publishes a zk-SNARK establishing that its update is exactly the
fixed-point gradient step computed from the current global model on a batch of
data the node committed to \emph{before} training, and in which the verifier ---
not the prover --- supplies the public inputs against which that proof is checked.
Nodes then agree on the \emph{aggregate} of the verified proposal set through a
leaderless two-phase Byzantine-fault-tolerant protocol: a quorum of matching
prevotes licenses a lock, a quorum of matching precommits forms a transferable
certificate, and any node holding one commits regardless of its own timers. Both
quorums are transferable objects rather than local observations, which is what
lets a node that missed them adopt the result instead of stalling. We
state the fault model explicitly ($n \geq 3f+1$, quorum $q > (n+f)/2$, partial
synchrony), prove agreement, and argue progress. Equivocation by either a voter
or a proposer yields self-verifying evidence that every honest node checks for
itself. Every committed transition is written to a hash-chained log that an
outside party can re-verify offline, re-checking each proof against the model as
of that height and recomputing each aggregate. Our circuit for a four-feature
logistic regression over a batch of eight samples compiles to 19{,}402 R1CS
constraints and proves in 0.84\,s; an ablation against a
6-constraint circuit of the kind used in prior proof-of-concept work shows the
latter admits a chosen-model forgery at the cost of one integer division per
parameter. Over 151 automated tests and five measured experiments we report
where the system holds and where it does not: coordinate-wise trimmed mean
loses 11.4\% accuracy precisely when fewer values are trimmed than
there are adversaries, and at 20\% packet loss the protocol makes
almost no progress at any parameter setting. We claim no privacy for the
published updates.
\keywords{Decentralized Federated Learning \and zk-SNARKs \and Byzantine Fault
Tolerance \and Consensus \and Verifiable Training \and Proof of Training}
\end{abstract}

% ─────────────────────────────────────────────────────────────────────────────
\section{Introduction}

Federated Learning (FL)~\cite{mcmahan2017fedavg} moves computation to the data
rather than the reverse. Removing the central aggregator removes a single point
of trust and failure, but it replaces one problem with a harder one: with no
trusted party, on what basis does a participant accept another participant's
model update?

The honest answer, in most decentralized FL systems, is \emph{none}. An update
is a vector of numbers. Nothing about the vector reveals whether it came from
running the agreed training algorithm on real data or from an optimizer pointed
at whatever objective the adversary preferred. Byzantine-robust aggregation
rules~\cite{blanchard2017krum,yin2018trimmedmean,cao2021fltrust} bound the
influence of a minority of such updates statistically, but they do not
\emph{detect} them, and their guarantees hold only under assumptions about how
many adversaries there are.

This paper takes the position that the update should carry its own proof, that
the proof should be checked against a statement the \emph{verifier} chooses, and
that the resulting protocol should state the fault model it needs rather than
assert resilience. Concretely we address three problems together.

\begin{enumerate}
  \item \textbf{An update that was never trained.} A node can publish any
    vector. We require a zk-SNARK that the published weights are the fixed-point
    gradient step from the current global model on a batch of samples whose
    Merkle root the node registered before training began. The data stays
    private; the fact that the computation was performed does not.
  \item \textbf{Disagreement about what was decided.} Nodes must commit the same
    global model at the same height even when a minority is actively hostile. We
    use a leaderless two-phase BFT protocol whose safety we prove under an
    explicitly stated fault model.
  \item \textbf{Misbehaviour that leaves no trace.} A Byzantine node that sends
    conflicting messages to different peers should not merely be suspected but
    \emph{provably} identified, by evidence any third party can check.
\end{enumerate}

\paragraph{An earlier design, and why it failed.}
This work began as a redesign of our own proof-of-concept, and the failures of
that prototype shape the contributions, so we state them plainly. Its circuit
had six R1CS constraints and accepted the gradient as a free witness signal: a
prover could choose the target weights, back-derive a gradient satisfying the
single subtraction, and obtain a valid proof. Soundness was vacuous --- every
update was provable. Its consensus committed the first valid update to arrive,
so the global model was a chain of single-node steps ordered by network latency
rather than an aggregate of anything. Its quorum, $n=10$ and $q=6$, was reported
as tolerating 30\% Byzantine nodes; $2 \cdot 6 > 10 + 3$ is false,
so it tolerates one. We report the difference between that design and this one
as a measured ablation rather than as a claim.

\paragraph{Contributions.}
\begin{itemize}
  \item A Circom circuit (Section~\ref{sec:circuit}) that binds an update to a
    committed dataset, a specific base model, and a specific learning rate, with
    every intermediate value range-checked and fixed-point division constrained
    by an explicitly range-checked remainder. 19{,}402 constraints for
    $n=4$ features and a batch of $k=8$.
  \item A verification discipline (Section~\ref{sec:verify}) in which the
    verifier constructs the public inputs from its own committed state and
    rejects any proof whose public signals differ, \emph{before} the pairing
    check.
  \item A leaderless two-phase consensus protocol (Section~\ref{sec:consensus})
    that decides the aggregate of the verified proposal set, with a stated fault
    model, a proof of agreement, and a liveness argument. We show that the
    single-phase alternative deadlocks on any network with non-zero delay, and
    that the two-phase one does too unless \emph{polkas} are gossiped as
    transferable certificates: without a proposer to carry a proof-of-lock, a
    split between locked and unlocked correct nodes is stable across every
    subsequent sub-round. Section~\ref{sec:eval-consensus} reports what the
    extra $n(n-1)$ deliveries per height cost.
  \item Accountable equivocation detection for \emph{both} voters and proposers,
    producing transferable evidence, together with the observation that
    detection alone is insufficient: an equivocating proposer's first message is
    accepted before its second exposes it, and failing to retract it splits the
    network permanently (Section~\ref{sec:accountability}).
  \item An evaluation (Section~\ref{sec:eval}) covering learning quality against
    three baselines, proof cost, consensus cost and message complexity, five
    adversarial scenarios, and offline auditability --- including the
    configurations in which the system fails.
\end{itemize}

All code, circuits, and the JSON result files backing every number in
Section~\ref{sec:eval} are provided as an anonymized artifact at
\texttt{[artifact URL --- to be inserted at submission]}. Each experiment is
regenerable with a single \texttt{make} target, and
\texttt{scripts/check\_paper\_numbers.py} re-checks every quantitative claim in
this paper against those files, so a stale number fails a script rather than a
reader.

% ─────────────────────────────────────────────────────────────────────────────
\section{Related Work}

\subsection{Threats to Federated Learning}

Mothukuri et al.~\cite{mothukuri2021survey} and Lyu et
al.~\cite{lyu2020threats} survey the attack surface: model poisoning, backdoors,
membership inference, and reconstruction. Kairouz et
al.~\cite{kairouz2021advances} identify verifying client behaviour without a
trusted third party as an open problem, which is the problem this paper attacks.

Two results bound what any scheme of our kind can claim about privacy. Zhu et
al.~\cite{zhu2019deepleakage} and Geiping et al.~\cite{geiping2020inverting}
reconstruct training samples from published gradients. Since every node in an FL
protocol needs the per-node update in order to aggregate it, \emph{publishing
the update is the leak}, and no amount of zero-knowledge in the proof changes
that. We return to this in Section~\ref{sec:limitations}.

\subsection{Zero-Knowledge Proofs of Training}

The work closest to our core claim proves that a model is the result of running
a training procedure. Garg et al.~\cite{garg2023zkpot} give the first practical
zero-knowledge proofs of training for logistic regression, and
Kaizen~\cite{abbaszadeh2024kaizen} extends this to deep networks with recursive
composition. Both operate at a far larger scale than we do and at
correspondingly greater cost.

We do not claim to improve on either. Our circuit proves a single fixed-point
gradient step on a small model; the contribution here is the \emph{composition}
of such a proof with a Byzantine-fault-tolerant agreement protocol and an
auditable log, and the demonstration that each individual piece must be built
much more carefully than a proof-of-concept suggests.

\subsection{zk-SNARKs Applied to Federated Learning}

zkFL~\cite{wang2024zkfl} has a central aggregator prove that it aggregated
correctly, which is the dual of our problem: it protects clients from the
server, whereas we protect participants from each other. ByzSFL~\cite{fan2025byzsfl}
combines zero-knowledge proofs with Byzantine-robust aggregation, and
Fu et al.~\cite{fu2025zkpconsensus} use ZK proofs inside a blockchain-secured FL
consensus. Federify~\cite{keshavarzkalhori2023federify} and
VeriBlock-FL~\cite{raghavarapu_veriblockfl} anchor verification on a blockchain.
These systems retain either a central aggregator or an external ordering
service.

\begin{table}[t]
\centering
\caption{Where this work sits. \textbf{Prover} is who produces the proof;
\textbf{Ordering} is what decides the sequence of global models;
\textbf{Decision} is what a height commits. A dash means the system does not
address that axis.}
\label{tab:related}
\begin{tabular}{llllc}
\toprule
\textbf{System} & \textbf{Prover} & \textbf{Ordering} & \textbf{Decision} &
\textbf{Fault model} \\
\midrule
zkFL~\cite{wang2024zkfl}        & aggregator & central server & server's sum & --- \\
Federify~\cite{keshavarzkalhori2023federify} & client & blockchain & on-chain agg. & chain's \\
VeriBlock-FL~\cite{raghavarapu_veriblockfl} & client & blockchain & on-chain agg. & chain's \\
VerifBFL~\cite{bellachia2025verifbfl} & client & blockchain & on-chain agg. & chain's \\
ByzSFL~\cite{fan2025byzsfl}     & aggregator & central server & robust agg. & robust-agg. \\
Byeon et al.~\cite{byeon2025trusted} & client & ring order & predecessor's model & --- \\
\midrule
This work & client & in-protocol BFT & aggregate of \emph{all} &
$n \ge 3f{+}1$, stated \\
\bottomrule
\end{tabular}
\end{table}

Table~\ref{tab:related} is the short version. The systems that verify training
either keep a central aggregator, in which case ordering is not a problem they
have, or delegate ordering to a blockchain, in which case the fault model is the
chain's and the latency and infrastructure come with it. The one leaderless
design in the list orders nodes in a ring, which makes each height depend on one
predecessor. What is missing from every row is a height that decides the
aggregate of \emph{all} verified proposals under a fault model the paper states
and enforces, which is the gap this work addresses.

\subsection{Decentralized FL}

Lalitha et al.~\cite{lalitha2019p2pfl} perform Bayesian model averaging over a
peer-to-peer graph with no aggregator but no verification either. Wang et
al.~\cite{wang2023enhancing} and Ferenczi and B\u{a}dic\u{a}~\cite{ferenczi2024decentralized}
coordinate decentralized FL through smart contracts, inheriting the blockchain's
latency and infrastructure requirements. Byeon et al.~\cite{byeon2025trusted}
arrange healthcare IoT nodes in a ring where each verifies its predecessor's
proof before forwarding; the sequential structure means one slow or absent node
delays the round. VerifBFL~\cite{bellachia2025verifbfl} combines verifiable
training with a blockchained FL pipeline.

Our design differs in that all $n$ nodes participate concurrently and
symmetrically, agreement is reached by an in-protocol BFT round rather than an
external chain, and the decision is the aggregate of \emph{all} verified
proposals rather than a selection among them.

\subsection{Byzantine-Robust Aggregation}

Krum~\cite{blanchard2017krum}, coordinate-wise trimmed mean and
median~\cite{yin2018trimmedmean}, and FLTrust~\cite{cao2021fltrust} bound the
influence of malicious updates on the aggregate. These are complementary to
proofs of training, not alternatives: a proof establishes that an update is the
correct gradient step on committed data, and an aggregation rule bounds the
damage when the committed data is itself adversarial. We use coordinate-wise
trimmed mean and report, in Section~\ref{sec:eval-adversarial}, exactly what
happens when its precondition is violated.

\subsection{Byzantine Consensus}

PBFT~\cite{castro1999pbft} established practical BFT in the partially
synchronous model~\cite{dwork1988partialsynchrony}, whose necessity follows from
FLP~\cite{fischer1985flp}. Tendermint~\cite{buchman2018tendermint} and
HotStuff~\cite{yin2019hotstuff} are the modern lineage. Our round structure is
Tendermint's, specialised to a leaderless setting where the value under
consideration is a deterministic function of the proposal set rather than a
proposer's choice; the safety argument in Section~\ref{sec:safety} is
Tendermint's argument adapted to that setting. Polygraph~\cite{civit2021polygraph}
introduced accountable Byzantine agreement, in which misbehaviour produces
irrefutable evidence; our equivocation handling is in that tradition.

% ─────────────────────────────────────────────────────────────────────────────
\section{System and Threat Model}
\label{sec:model}

\subsection{Participants}

A permissioned set of $n$ nodes, fixed for a run and known to all. Each node
holds an Ed25519 key pair; its identity is the hex encoding of its public key.
Before height~1, each node registers a Merkle root $\rho_i$ over its quantized,
batched dataset together with a signature; the binding from node to root is
immutable thereafter.

This provides \textbf{no Sybil resistance}. An adversary able to enrol
identities at will controls the vote and the model, and no amount of proving
changes that. Admission control is out of scope and must be provided by whatever
operates the deployment.

\subsection{Fault Model}

At most $f$ nodes are Byzantine: arbitrary behaviour, including equivocation,
collusion among all $f$, adaptive timing, and full visibility of the network.
The remaining $n-f$ nodes are correct, though they may be slow, may crash and
restart (recovering from an on-disk log), and may be temporarily partitioned.
Communication is \emph{partially synchronous}~\cite{dwork1988partialsynchrony}:
messages between correct nodes are eventually delivered, and after an unknown
global stabilisation time delays are bounded by an unknown constant $\Delta$.
Byzantine nodes cannot forge signatures and cannot break the discrete-logarithm
assumptions on \textbf{bn128}.

The quorum $q$ must satisfy
\begin{equation}
  2q > n + f \quad \text{(safety)}, \qquad n - f \geq q \quad \text{(liveness)},
  \label{eq:quorum}
\end{equation}
which together force $n \geq 3f+1$. The first inequality makes any two quorums
share at least $2q-n > f$ nodes, hence at least one correct node; the second
makes the correct nodes alone able to form a quorum. Our implementation refuses
to start with any $(n, f, q)$ violating~(\ref{eq:quorum}), rather than allowing a
run to advertise a guarantee it does not hold. As noted in the introduction, the
parameters $n=10$, $q=6$ used in our earlier prototype satisfy~(\ref{eq:quorum})
only for $f=1$.

\subsection{What the System Does Not Defend}

Stated here so that no later claim is read past its scope; expanded in
Section~\ref{sec:limitations}. Updates are public and therefore invertible.
Groth16 requires a per-circuit trusted setup. A node that commits to a poisoned
dataset and trains on it correctly produces a valid proof, and nothing binds a
batch to a height, so it may keep re-training on its worst one. Trimmed-mean
robustness is conditional on trimming at least $f$ values per side. Sustained
packet loss defeats the protocol. Nothing here is constant-time.

% ─────────────────────────────────────────────────────────────────────────────
\section{The Verifiable Update Circuit}
\label{sec:circuit}

\subsection{Fixed-Point Representation}

zk-SNARK arithmetic is over a prime field, so we represent a real value $v$ as
the integer $\tilde v = \mathrm{round}(v \cdot S)$ with $S = 10^6$, negative
values being encoded as field elements $p - |\tilde v|$. Quantization on the
Python side uses exact rational arithmetic with ties resolved away from zero, so
that trainer and circuit agree bit-for-bit rather than approximately.

Every signal is range-checked with \texttt{Num2Bits}: features satisfy
$|\tilde x| < 2^{25}$, weights and bias $|\tilde w|, |\tilde b| < 2^{26}$, and
the pre-activation $|\tilde z| < 2^{35}$, where the last bound is
\emph{derived} from the first two and the feature count rather than chosen. This
matters: without range checks a prover can present a huge field element that the
circuit's own arithmetic then reads as a small negative number, and the
constraints are satisfied by a witness that means nothing.

\subsection{Fixed-Point Division}
\label{sec:div}

The single subtlest component. Circom's \verb|\| and \verb|%| operate on field
elements, so for a negative numerator they divide $p - |{\rm num}|$ and return a
quotient that is arithmetically meaningless while still satisfying a naive
constraint. Our \texttt{DivConst} template computes Euclidean division that
matches Python's \verb|//| exactly, including for negative numerators, by
offsetting before dividing:

\begin{lstlisting}
var SHIFT = divisor * (1 << quotBits);
q <-- (num + SHIFT) \ divisor - (1 << quotBits);
r <-- (num + SHIFT) % divisor;
num === q * divisor + r;
Num2Bits(remBits)(r);            // 0 <= r
Num2Bits(remBits)(divisor-1-r);  // r < divisor
SignedRangeCheck(quotBits)(q);   // q in range
\end{lstlisting}

The two range checks on $r$ are what make the constraint
$\mathrm{num} = q\,d + r$ pin down a unique $(q, r)$; without them any $q$ admits
a compensating $r$ and the division constrains nothing. Our earlier circuit's
\verb|dot_product * SCALE === sum| had exactly this defect.

\paragraph{The numerator needs a bound too.}
$\mathrm{num} = q\,d + r$ is an equation \emph{in the field}, so the checks on
$q$ and $r$ isolate a unique integer pair only while the numerator cannot wrap:
we need
\begin{equation}
  |\mathrm{num}| \;\le\; d \cdot 2^{\rm quotBits} ,
  \label{eq:numbound}
\end{equation}
or a large field element satisfies the equation with an in-range $(q, r)$ that
means nothing. This is the same class of defect as the missing remainder check
and deserves the same treatment, so we state the bound for every division in the
circuit rather than leaving it to be re-derived.

Each numerator is an \emph{assigned} combination of range-checked signals
(\texttt{<==}), never a free witness, so each has a derivable bound: the forward
pass divides $\langle \tilde{\mathbf{x}}, \tilde{\mathbf{w}}\rangle$, at most
$n 2^{x+w} = 2^{53}$, by $S$ with a 35-bit quotient, giving $2^{55}$ of room;
the weight update divides $\eta \sum_t e^{(t)} \tilde x^{(t)}_j$, at most
$2^{68}$, by $kS^2$ with a 32-bit quotient, giving $2^{75}$; the bias update
needs $2^{43}$ against $2^{55}$. \texttt{FixedPointConfig.check()} re-derives
all four at start-up and refuses a configuration that violates
(\ref{eq:numbound}), so widening a feature vector cannot silently invalidate the
argument.

The sigmoid was the interesting case, and writing this check is what found it.
Its numerator is $m_{\rm sel}(a - t_{\rm sel})$, which is \emph{not} bounded by
$|z|$: on the flat outer piece the slope $m_{\rm sel}$ is zero, so a large $|z|$
contributes nothing, and elsewhere $|a - t_{\rm sel}|$ is at most the last
breakpoint. That bound held for the honest witness --- the one-hot selector is
determined by $a$ --- but it was \emph{derived from the comparator chain being
sound} rather than enforced anywhere, and (\ref{eq:numbound}) has to hold for
every witness, not the honest one. We now constrain the product directly with an
\texttt{AbsBound}, at $1.16 \times 10^{12}$ against a limit of $1.68 \times
10^{13}$. Together with an explicit range check on the activation output ---
which the template advertised as lying in $[0, S]$ and did not enforce, though
every downstream numerator bound depends on it --- this costs 1{,}048
constraints, about 5.7\%, and replaces a chain of reasoning with a check.

\subsection{Piecewise-Linear Sigmoid}

A faithful sigmoid is expensive in R1CS. We use a five-piece linear
approximation with breakpoints at $|z| \in \{1, 2, 3, 5\}$, reflected about
$(0, \tfrac{1}{2})$, selected by four comparators feeding a one-hot selector
built from monotone differences, and evaluated with one multiplication and one
\texttt{DivConst}. Its maximum absolute error against the true sigmoid is
$0.0117$, measured over the representable domain. This is an approximation, and
Section~\ref{sec:limitations} states what it costs.

\subsection{Statement Proved}

The template $\textsf{ModelUpdateCheck}(n, k, d, S, \cdot)$ takes public
$(\mathbf{w}_{\rm old}, b_{\rm old}, \mathbf{w}_{\rm new}, b_{\rm new}, \eta,
\rho)$ and private $(\mathbf{x}^{(1..k)}, y^{(1..k)})$ together with $d$-level
Merkle paths, and enforces, in order:

\begin{enumerate}
  \item range checks on every input signal;
  \item $y^{(t)}(y^{(t)} - 1) = 0$ for each sample, and $0 \le \eta \le S$;
  \item \textbf{Merkle inclusion} of the batch digest against the public root
    $\rho$, where the batch digest is a Poseidon~\cite{grassi2021poseidon} chain
    over per-sample digests, so that the batch is bound to the committed dataset
    \emph{and} to its ordering;
  \item the forward pass $\tilde z^{(t)} = \langle \tilde{\mathbf{x}}^{(t)},
    \tilde{\mathbf{w}}_{\rm old}\rangle / S + \tilde b_{\rm old}$, the
    piecewise-linear activation, and the error
    $e^{(t)} = \widetilde{\rm pred}^{(t)} - y^{(t)} S$;
  \item the parameter update, by \emph{constraint} rather than computation:
    \begin{equation}
      \tilde w_{{\rm new}, j} \;=\; \tilde w_{{\rm old}, j}
        \;-\; \left\lfloor \frac{\eta \sum_t e^{(t)} \tilde x^{(t)}_j}
                               {k\,S^2} \right\rfloor,
      \qquad
      \tilde b_{\rm new} \;=\; \tilde b_{\rm old}
        - \left\lfloor \frac{\eta \sum_t e^{(t)}}{k\,S} \right\rfloor
      \label{eq:update}
    \end{equation}
    each floor division being a \texttt{DivConst} with a range-checked
    remainder;
  \item a verifiable clipping bound $|\tilde w_{{\rm new}, j} - \tilde w_{{\rm
    old}, j}| \le \delta_{\max}$ and $|\tilde b_{\rm new} - \tilde b_{\rm old}|
    \le \delta_{\max}$, enforced with two \texttt{Num2Bits} rather than a
    power-of-two range check, so the bound is exactly $\delta_{\max} = 5 \times
    10^5$ and not the next power of two above it.
\end{enumerate}

Note that $b_{\rm new}$ is constrained. In the earlier circuit it was not, which
left the bias entirely free.

\subsection{Verification: the Statement Is the Verifier's}
\label{sec:verify}

A proof is meaningless if the verifier checks it against public inputs the
prover supplied --- that establishes only that the prover was self-consistent. In
our system the verifier constructs
$(\mathbf{w}_{\rm old}, b_{\rm old}, \eta, \rho_{\rm proposer})$ from
\emph{its own} committed model, its own configured learning rate, and the
proposer's registered root, and rejects any proof whose public signals differ
from that expectation \emph{before} performing the pairing check. The prover
supplies the proof and $(\mathbf{w}_{\rm new}, b_{\rm new})$; it supplies
nothing else.

The ordering of public signals is a convention shared between the Circom
template and the host code. We check it against the compiled \texttt{.sym} file
at setup time rather than trusting the convention, since a silent reordering
would make the verifier check a different claim than it believes it is checking.

\subsection{Trusted Setup}

Groth16~\cite{groth2016size} requires a per-circuit ceremony, and whoever holds
its toxic waste can forge proofs for that circuit. Our setup script runs a real
multi-contributor phase-2 ceremony with a final beacon over a public
powers-of-tau transcript, which is sound as long as one contributor discards
their randomness --- a natural fit for a peer-to-peer system, where the
participants are already present. An offline fallback generates phase~1 locally;
in that mode the operator can forge proofs, the script says so, and the build
metadata records it. We regard a paper whose premise is ``no trusted party''
as obliged to state this rather than omit it.

% ─────────────────────────────────────────────────────────────────────────────
\section{Consensus}
\label{sec:consensus}

\subsection{Aggregate, Do Not Select}

A height decides $\mathrm{Agg}(\mu_{h-1}, P)$ over the entire verified proposal
set $P$, not one chosen element of it. This is the difference between federated
learning and a race: committing the first valid update to arrive, as our earlier
design did, makes the global model a chain of single-node steps ordered by
network latency.

It has a second consequence that the protocol depends on. Because $\mathrm{Agg}$
is a deterministic function of public inputs, two correct nodes holding the same
proposal set necessarily vote identically. A conflicting vote is therefore
\emph{misbehaviour} rather than a timing artefact, which is precisely what makes
an accusation of equivocation meaningful.

$\mathrm{Agg}$ orders $P$ canonically by proposal identifier, applies a
coordinate-wise rule --- trimmed mean, median, or weighted mean --- and performs
all arithmetic in integers with floor division. Every step is reproducible by
any observer from public data. The trimmed mean drops
\begin{equation}
  t \;=\; \min\!\left(\left\lfloor \frac{\pi\,|P|}{1000} \right\rfloor,\;
                      \left\lfloor \frac{|P|-1}{2} \right\rfloor \right),
  \qquad \pi = 250,
  \label{eq:trim}
\end{equation}
values from each end of every coordinate. The second term of the minimum is not
cosmetic: it keeps at least one value, so the rule is total on any non-empty
proposal set rather than raising on small ones. It also means the robustness
precondition $t \geq f$ \emph{fails silently for small $|P|$} --- at $|P|=4$ and
$\pi=250$, $t=1$, so a run with $f=3$ is trimming far less than its own fault
model assumes. We therefore compute $t$ and $f$ at every height and log the
comparison rather than leaving it implicit; Section~\ref{sec:eval-adversarial}
reports a run where the bound is violated and the aggregate degrades exactly as
the arithmetic predicts.

\subsection{A Height}

Each height proceeds in sub-rounds, each sub-round in two phases:

\begin{description}
  \item[Collect.] Proposals arrive and are verified (Section~\ref{sec:verify})
    before being stored or relayed. Nothing unverified ever leaves a node.
  \item[Prevote.] On timeout, a node broadcasts $\mathrm{prevote}(h, r,
    H(\mathrm{Agg}(\mu_{h-1}, P)))$ --- or, if it holds a lock, its locked value.
  \item[Polka and lock.] $q$ matching prevotes at sub-round $r$ constitute a
    \emph{polka}: proof that a quorum was willing to commit that value. Only a
    polka licenses a node to lock. A polka is a first-class transferable
    certificate, gossiped like a commit certificate and verified from its $q$
    signatures by whoever receives it --- without that, a node which merely
    missed the prevotes can never leave a split (Section~\ref{sec:liveness}).
  \item[Precommit.] A node that locked broadcasts $\mathrm{precommit}(h, r, v)$.
    A node that saw no polka before its timer expires precommits \textsc{nil}
    and moves to sub-round $r+1$ with a doubled collect window.
  \item[Commit.] $q$ matching precommits form a \emph{certificate}. Any node
    holding one commits the height, appends to its audit log, and advances ---
    regardless of its own timers.
\end{description}

\paragraph{A height with no proposals.}
A node whose verified proposal set is empty has nothing to aggregate and so
prevotes nothing at all: it opens the next sub-round with a longer window rather
than prevoting \textsc{nil}, since \textsc{nil} is the outcome of a failed
polka, not a statement about an empty set. The height simply does not advance
until some proposal arrives, which is why Proposition~2 assumes at least one
correct node proposes. This is a liveness dependency and not a safety one: with
no proposals there is no value to certify, and nothing can be committed by
anyone. Proposing is decoupled from voting --- every node votes at every height,
but only nodes holding data propose --- so a deployment in which some members
are verifiers only is normal, and one in which \emph{no} correct member holds
data does not make progress.

\subsection{Why Two Phases}

The two phases are not ceremony, and the single-phase alternative is not a
simplification but a deadlock. Without a polka, a node has nothing to lock on
except its own prevote. Correct nodes close their collect windows with different
proposal sets on any network with non-zero delay, so they compute different
aggregates, lock on different values, and no value ever reaches $q$ again. We
observed exactly this in our own implementation: two nodes, two locks, zero
rejections, no certificate, indefinitely.

Locking on a polka instead means that split prevotes produce \emph{no} lock:
every node precommits \textsc{nil}, sub-round $r+1$ opens with a longer window,
and by then the proposal sets have converged. This is Tendermint's
rule~\cite{buchman2018tendermint}. Converged proposal sets are enough only while
\emph{no} node holds a lock, since a locked node prevotes its lock rather than
its aggregate; Section~\ref{sec:liveness} shows that once one does, the two
phases are still not sufficient unless the polka itself can be forwarded.

\subsection{Safety}
\label{sec:safety}

\begin{proposition}[Agreement]
If $2q > n+f$ and at most $f$ nodes are Byzantine, no two correct nodes commit
different models at the same height.
\end{proposition}

\begin{proof}
Write $v \neq v'$ for the two committed hashes and $r \le r'$ for the sub-rounds
of their certificates $C$ and $C'$. Two quorums intersect in at least $2q - n >
f$ members, so every pair of quorums contains a common \emph{correct} node.
Note also that $n \ge 3f+1$ and $2q > n+f$ give $q \ge 2f+1$, so each certificate
contains at least $q - f \ge f+1$ correct precommits.

The argument rests on one property of the implementation, which we state
separately because it is what makes the induction close:

\begin{quote}
\textbf{(I)} A correct node never prevotes at sub-round $s$ after observing a
polka at $s$. Observing a polka sets its step to \textsc{precommit}, and the
prevote action fires only in the \textsc{prevote} step.
\end{quote}

Consequently a correct node's prevote at $s$ reflects a lock formed
\emph{strictly before} $s$. Because the proof leans on it, \textbf{(I)} is a
test rather than a remark: a node is handed a polka before its own prevote timer
fires, and the test fails if any prevote for that sub-round is ever emitted.

\medskip\noindent\emph{Case $r = r'$.} $C$ and $C'$ are quorums of precommits in
the same sub-round, so some correct node signed both, i.e.\ two different votes
in the slot $(p, h, r, \textsc{precommit})$. A correct node signs one vote per
slot. Contradiction.

\medskip\noindent\emph{Case $r < r'$.} We show the stronger claim, stated over
the \emph{certificate} rather than over a single vote, because that is what the
induction step consumes:

\begin{quote}
\textbf{(C)} If a commit certificate for $v$ exists at sub-round $r$, then for
every $r_1 \ge r$, any polka at $r_1$ is for $v$.
\end{quote}

Since $C'$ exists, some correct node precommitted $v'$ at $r'$, which it does
only immediately after locking on $v'$ --- so a polka for $v' \neq v$ exists at
$r' \ge r$, and \textbf{(C)} applied to $C$ gives $v' = v$, a contradiction. It
remains to prove \textbf{(C)}, by strong induction on $r_1$.

\emph{Base $r_1 = r$.} $C$ contains at least $q - f \ge f+1$ correct precommits
for $v$ at $r$, and a correct node precommits a non-\textsc{nil} value only
after observing a polka for it, so a polka for $v$ at $r$ exists. Suppose a
polka for $u$ at $r$ also exists. The two polkas are quorums of prevotes in the
same sub-round, so some correct node signed both, i.e.\ two votes in the slot
$(p, h, r, \textsc{prevote})$. Hence $u = v$.

\emph{Step.} Let $r_1 > r$ and assume \textbf{(C)} for every sub-round in $[r,
r_1)$. Suppose a polka for $u$ at $r_1$. Its $q$ prevotes and the $q$ precommits
of $C$ --- both quorums, which is why the claim is stated over $C$ and not over
one node's vote --- share a correct node $p^*$, which therefore precommitted $v$ at $r$ and
prevoted $u$ at $r_1$. Precommitting $v$ at $r$ means $p^*$ held the lock $(v,
r)$. A correct node prevotes its locked value, so if $u \neq v$ then $p^*$'s lock
moved off $v$ before it prevoted at $r_1$; locks move only on an observed polka,
and only at a sub-round no earlier than the locked round. So $p^*$ observed a
polka for some $u_0 \neq v$ at a sub-round $r_0 \in [r, r_1]$, before prevoting
at $r_1$. By \textbf{(I)}, $r_0 \neq r_1$: a node that has seen a polka at $r_1$
does not go on to prevote at $r_1$. Hence $r_0 \in [r, r_1)$, and the induction
hypothesis makes that polka a polka for $v$, i.e.\ $u_0 = v$ --- contradicting
$u_0 \neq v$. Therefore $u = v$.
\qed
\end{proof}

Since entry $h$ records $H(\mu_{h-1})$ and the audit chain hashes each entry
together with its predecessor's link, two correct nodes agreeing at every height
up to $h$ hold bit-identical logs. Our simulator asserts this after every run.

Proposition~1 says nothing about the \emph{quality} of $\mu_h$: $f$ Byzantine
nodes can submit valid proofs for genuinely-trained but adversarially-chosen
updates, and only the aggregation rule bounds their influence. It also assumes
exactly the model of Section~\ref{sec:model}; with $f+1$ Byzantine nodes two
conflicting certificates are constructible and the argument fails at its first
step, as it must.

\subsection{Liveness}
\label{sec:liveness}

\begin{proposition}[Progress]
Under partial synchrony, with $n - f \ge q$, at most $f$ Byzantine nodes, and at
least one correct node proposing, every height eventually commits at some
correct node.
\end{proposition}

\begin{proof}[Sketch]
After stabilisation, delays are bounded by $\Delta$. The
collect window is $w \cdot 2^r$, so some sub-round $r^*$ has $w 2^{r^*} >
2\Delta$ --- which presumes sub-rounds continue without bound, a hypothesis our
implementation violated for some time (see below). From $r^*$ onward every
message sent at the start of a phase arrives before that phase closes at every
correct node.

It does \emph{not} follow that correct nodes prevote identically merely because
they hold the same proposal set, and this is the step at which an earlier
version of both our argument and our implementation was wrong. A locked node
prevotes its \emph{lock}, not $\mathrm{Agg}(P)$. If some correct nodes locked on
$v$ at an earlier sub-round and the rest never observed that polka, the locked
ones offer $v$ and the unlocked ones offer $\mathrm{Agg}(P) = v' \neq v$ once
further proposals have arrived. Should neither camp reach $q$, no polka forms;
but only a polka moves a lock, so the split reproduces itself at every
subsequent sub-round and the height never finishes. Extra sub-rounds do not
help, because each one recreates the same two camps.

Transferable polkas are what close this. A node that observes $q$ matching
prevotes gossips them as a certificate; a node receiving one verifies the $q$
signatures itself and adopts the lock whenever the polka's sub-round is at least
its own locked round. After $r^*$ every such certificate reaches every correct
node, so all correct nodes hold the lock of the highest sub-round any of them
has seen, and prevote it. That is $n - f \ge q$ identical prevotes, hence a
polka, hence $q$ matching precommits, hence a certificate. Byzantine nodes can
withhold or send garbage, neither of which prevents $q$ correct nodes from
agreeing. \qed
\end{proof}

This is the leaderless counterpart of Tendermint's proposer re-proposing with a
proof-of-lock from the highest round~\cite{buchman2018tendermint}. With no
proposer to carry it, the evidence has to travel by itself. It costs one extra
gossiped message per polka, which Section~\ref{sec:eval-consensus} accounts for
separately rather than folding into the vote total.

Four implementation details are load-bearing here, and each was a defect in our
implementation at some point. The window must \emph{grow}, or no $r^*$ exists ---
a run on loopback, where $\Delta \approx 0$, hides this entirely. A quorum of
\textsc{nil} precommits must start the next sub-round immediately rather than
waiting out a timer whose outcome overlapping quorums have already settled. And
a node that falls behind must \emph{keep} asking for the missing height: asking
once and never retrying leaves it stuck forever whenever the answer is lost.

The fourth is the one the proposition's proof makes obvious in hindsight, and it
is the most severe. The proof needs sub-rounds to keep opening; our
implementation capped them at a parameter $r_{\max}$, past which a node asked a
peer to send it the height instead. But a sync request can only be answered by a
node that is \emph{ahead}, and the situation that exhausts the budget is
precisely the one in which no quorum formed --- so nobody is ahead, nothing
answers, and because every correct node reaches the budget at roughly the same
time, all of them fall silent together. The implementation contradicted its own
liveness argument, and the code was not a faithful rendering of the protocol we
had proved. A node now keeps its clock running and asks for sync alongside;
$r_{\max}$ is the point at which it starts asking, not a ceiling. We report the
consequence in Section~\ref{sec:eval-consensus}, because it accounts for
behaviour we had been reporting as a limit of the design.

\paragraph{A correct node must not forget that it voted.}
Both cases of Proposition~1 rest on \emph{a correct node signs one vote per
slot}, and the fault model admits crash-restart. A node that crashed after
broadcasting a prevote and restarted without it would sign a second, different
vote in the same slot: not Byzantine, but indistinguishable from it, blacklisted
by its peers on evidence that is entirely genuine, and a counterexample to the
hypothesis the proof uses. So a vote is appended to a write-ahead log and
\texttt{fsync}ed \emph{before} it is broadcast, and on restart a slot already in
the log re-sends exactly the vote it recorded. A crash between recording and
sending loses the broadcast, which is harmless --- peers time out and the
sub-round proceeds --- rather than losing the record, which is not. Simulations
never restart and use an in-memory log that says so; a deployment passes a path.

\paragraph{Delay and loss are different problems.}
The doubling window answers delay. It does not answer loss: a dropped vote costs
the whole sub-round however long the window was, and a two-phase sub-round needs
two quorums to land rather than one. What loss needs is more \emph{attempts},
capped by a parameter $r_{\max}$. But the two are coupled --- the $r$-th attempt
uses a window of $w 2^{r}$, so once that exceeds the time a node is willing to
spend on a height, further attempts are unreachable and $r_{\max}$ stops
mattering. Section~\ref{sec:eval-consensus} reports what we measured, which is
less tidy than this reasoning suggests.

\subsection{Accountability}
\label{sec:accountability}

A voter's slot is $(\mathrm{voter}, h, r, \mathrm{phase})$; occupying one twice
is equivocation. Prevoting and precommitting within one sub-round is normal,
which is why the phase belongs in the slot. A proposer's slot is
$(\mathrm{proposer}, h)$.

The two conflicting signed messages \emph{are} the evidence. A recipient
re-verifies both signatures itself rather than trusting the accuser, so evidence
is transferable in the sense of Polygraph~\cite{civit2021polygraph}. Detection
happens on the signed message, before any expensive verification, so a proposer
flooding conflicting proposals is caught rather than paid for in proof
verifications.

\paragraph{Detection is not enough.}
An equivocating proposer's first message is accepted before its second exposes
it, and honest nodes observe the two in different orders --- so they accept
\emph{different} firsts. If the node that catches the equivocation locally keeps
what it already accepted while the nodes that learn by gossip discard it, the
proposal sets never converge. In a 10-node run this produced prevotes split 5/4
against $q = 7$ at every sub-round, with the offender correctly blacklisted
everywhere and not a single height committed. The fix is to retract: a proven
equivocator's votes \emph{and} proposals are erased from every open round and
from the future buffer. The purge is a deterministic function of self-verifying
gossiped evidence, so every honest node performs the same one.

\paragraph{What purging may and may not undo.}
Purging retroactively changes who counts, so its interaction with
Section~\ref{sec:safety} has to be pinned down rather than left implicit. We fix
the semantics as follows.

\begin{quote}
\textbf{(P)} A certificate or polka is valid \emph{once formed}. Validity is a
property of the $q$ signatures it carries and of the registered member set ---
never of which equivocations the checking node happens to have heard about.
Evidence governs \emph{future} tallies only: a blacklisted node's new votes are
refused and its pending votes and proposals are erased from open rounds, but
nothing already certified is unmade.
\end{quote}

Without \textbf{(P)} the proof does not apply, because ``a polka at $r_1$''
would not be a well-defined object: two correct nodes holding different evidence
would disagree about whether the same $q$ signatures constituted one, and a lock
licensed at one node would be unlicensed at another. With \textbf{(P)} every
predicate the proof quantifies over is a function of signed messages alone, so
all correct nodes agree on it, and the argument goes through unchanged.

\textbf{(P)} also costs nothing in safety. The whole point of $2q > n+f$ is to
budget for all $f$ faulty nodes signing both sides of a disagreement; an
equivocator's signature inside a quorum is precisely the case the inequality is
sized for, and refusing such a quorum neither prevents a conflicting certificate
nor is needed to. Our implementation originally did refuse them, which was
view-dependence introduced for no gain, and a test now pins the corrected
behaviour.

Locks, correspondingly, are not released: a lock is licensed by a polka, a polka
is valid once formed by \textbf{(P)}, and so the lock remains licensed no matter
whose votes are later purged. The next polka moves it by the ordinary rule.

This is accountability, not prevention: $f$ equivocators can still deny a height
and force additional sub-rounds. What they cannot do is deny it without leaving
a proof that they did.

\subsection{Auditability and Catch-Up}

Each committed height writes an entry containing the height, the parent model
hash, the new model and its hash, the proposal identifiers, the proposals with
their proofs, the registered roots they were checked against, the certificate,
and the aggregation parameters. Entries are hash-chained.

The same entry serves three purposes. It is the sync protocol: a node at height
$h$ requests entry $h+1$, verifies it end to end --- certificate quorum, every
proof against the model as of that height, recomputed aggregate, chain link ---
and applies it. It is the bootstrap path for a new node, from genesis. And it is
what an outside auditor re-verifies offline, sharing no state with the node that
produced it.

% ─────────────────────────────────────────────────────────────────────────────
\section{Implementation}

Python~3.12 with asyncio; circuits in Circom~2.1.6 compiled with
circom~2.2.3 and snarkjs~0.7.6 over bn128; Ed25519 signatures.

The consensus engine is a pure state machine: messages in, \emph{effects} out
(\textsf{Broadcast}, \textsf{Committed}, \textsf{NeedSync}, \textsf{OpenWindow},
\textsf{Rejected}). It performs no I/O. This is what makes the protocol testable
without sockets or sleeps, and it is why the tests in
Section~\ref{sec:eval-tests} are protocol tests rather than timing tests.

Four details are worth recording because each was a source of real bugs.
\emph{Canonical encoding:} all signed and hash-chained content is canonical JSON
that \emph{rejects floats outright}, since a float's shortest representation is
platform-dependent and a signature over a platform-dependent encoding is not a
signature; timestamps are integer milliseconds for the same reason.
\emph{Domain separation:} signatures cover
\texttt{"zkfl/v2" \textbar{} kind \textbar{} canonical}, so a signature over one message type
cannot be replayed as another. \emph{Proposal identity:} a proposal's identifier
hashes the entire payload including height and proposer, because an identifier
derived from the update alone collides whenever two nodes take the same step,
and then two distinct proposals share one identity. \emph{Vote durability:} a
node's own votes go to an \texttt{fsync}ed write-ahead log before broadcast, and
a restart re-sends what the log holds rather than signing again --- without it a
crash-restart is self-equivocation and Proposition~1's hypothesis is false.

Four message kinds cross the wire: proposals, votes, polkas and commit
certificates. The last three are self-verifying and terminal, so they are
gossiped; votes are sent once by their author and not relayed, which is what
keeps the cost quadratic rather than cubic.

A persistent Node.js worker holds the proving key in memory across proofs;
proving runs in a thread so a node that is proving continues to serve messages
rather than appearing partitioned. The transport is either an in-process bus
with injected one-way delay, jitter and loss, or real WebSockets; both run the
identical node code. A message gate applies per-peer rate limits and a bounded
future-height buffer before any verification, so flooding is paid for in
bookkeeping rather than in pairing checks.

% ─────────────────────────────────────────────────────────────────────────────
\section{Evaluation}
\label{sec:eval}

All measurements on an Intel Core i5-11400H at 2.70\,GHz, Node
v24.20.0. Each experiment writes a JSON file recording its numbers alongside the
CPU and toolchain versions, and each is regenerable with one command.

The three kinds of number in this section are not equally stable, and we
separate them rather than presenting them as one class of measurement.
Constraint counts, proof sizes, accuracies and tampering outcomes are
\emph{exact}: they are the same on every run. Message counts are exact only so
long as no collect window expires --- at $n \le 20$ every height is identical
down to the delivery, while at $n = 50$ a height takes tens of seconds and a
loaded host can expire a window and add a retried sub-round. Proof and
verification timings move a few percent. The consensus figures under packet loss
move far more than that, for reasons we quantify in
Section~\ref{sec:eval-consensus} and do not want a reader to discover by failing
to reproduce them.

\subsection{Test Suite}
\label{sec:eval-tests}

151 automated tests, of which 15 require the compiled circuit. They cover
fixed-point and Poseidon agreement with the circuit, in-circuit soundness,
public-input binding, quorum arithmetic, locking and certificates, equivocation
evidence, aggregation determinism, audit re-verification and tampering, message
gating, and live networks with each adversary in turn. We quote the count from a
run rather than from memory; our earlier draft reported 44 in its text and 53 in
its table.

\subsection{Circuit and Proof Cost}
\label{sec:eval-proof}

\begin{table}[t]
\centering
\caption{Proof cost, and an ablation against the earlier circuit. Cold includes
Node start-up and proving-key load; warm reuses the persistent worker. The
ablation's timings ($\dagger$) are \emph{not} comparable to this work's: that
circuit is driven through the \texttt{snarkjs} command line, so every figure in
its row pays a fresh Node process, which is why a 6-constraint circuit appears
to verify more slowly than a 19{,}402-constraint one. Constraints and samples
are the columns to read across rows.}
\label{tab:proof}
\begin{tabular}{lrrrrrr}
\toprule
\textbf{Circuit} & \textbf{Constraints} & \textbf{Samples} &
\textbf{Prove cold} & \textbf{Prove warm} & \textbf{Verify} & \textbf{zkey} \\
\midrule
This work ($n{=}4,k{=}8,d{=}4$) & 19{,}402 & 8 & 1.09\,s &
  0.84\,s & 15.1\,ms & 9.5\,MB \\
Earlier prototype (ablation)    & 6        & 0 & 0.52\,s$^\dagger$ &
  0.41\,s$^\dagger$ & 336.9\,ms$^\dagger$ & $<$0.1\,MB \\
\bottomrule
\end{tabular}
\end{table}

Table~\ref{tab:proof} gives the cost. Verification is 15.1\,ms against a
persistent worker. The 742\,B proof size is \texttt{snarkjs}' JSON encoding,
which is what our nodes actually put on the wire; the underlying Groth16 proof
is three \textbf{bn128} group elements, 128\,B compressed, and a deployment that
cared about bandwidth would send that instead. We quote the encoding we
measured rather than the theoretical minimum, and the message-cost table in
Section~\ref{sec:eval-consensus} is computed from the former.

The ablation row is the more informative one, and the column to read is
\emph{constraints}, not time. Our earlier prototype proved in
0.41\,s --- but on a harness that spawns a process per operation, so its timings
say more about Node than about cryptography, and we mark them accordingly rather
than inviting a speedup claim we did not measure. What the row does establish is
what the two circuits prove: reporting 0.41\,s as the cost of proving a training
step compares a proof of nothing with a proof of something. Concretely, the
\textsc{poison1} attack --- choose a target model, then back-derive the gradient
witness that satisfies the single subtraction constraint --- is accepted by that
circuit at a cost of one integer division per parameter and requires no training
data at all. Against the circuit of Section~\ref{sec:circuit} it fails at the
clipping-bound constraint; a variant that trains genuinely but on uncommitted
data fails at the Merkle inclusion constraint. Both failures are witness
generation errors, i.e.\ no proof exists to submit.

\subsection{Learning}
\label{sec:eval-learning}

\begin{table}[t]
\centering
\caption{Iris (versicolor vs.\ virginica), 100 samples, 4 features, 70/30 split,
5 nodes, 60 heights, trimmed mean, $\eta = 0.5$. Accuracy on the 30-sample test
set. \textbf{Range} is across 3 splits (seeds 0--2); the single-split column is
seed 0, the one the loss column is measured on.}
\label{tab:learning}
\begin{tabular}{lccc}
\toprule
\textbf{System} & \textbf{Accuracy} & \textbf{Range, 3 splits} &
\textbf{Final loss} \\
\midrule
This work (integer fixed-point, PWL sigmoid, BFT) & 0.867 & 0.867--0.967 & 0.353 \\
FedAvg, same partition, no proofs & 0.833 & 0.833--0.967 & 0.306 \\
Centralized float64 logistic regression & 0.833 & 0.833--0.967 & 0.269 \\
scikit-learn, same split & 0.833 & 0.833--0.967 & --- \\
\bottomrule
\end{tabular}
\end{table}

All 60 heights committed, all nodes bitwise identical, every invariant held.

The number to take from Table~\ref{tab:learning} is not $0.867 > 0.833$. The
test set holds 30 samples, so that gap is \emph{one sample} and is noise, and
the two further splits say so directly: on both of them this work and FedAvg
score \emph{identically} (0.933 and 0.967), and the spread across splits ---
0.100 for this work, 0.133 for FedAvg --- is three to four times the gap being
compared. We report three splits rather than one for exactly this reason. The
claim the table supports is the weaker and more useful one: quantizing to integers,
replacing the sigmoid with five line segments, clipping every step to
$\delta_{\max}$, and routing every update through a BFT protocol does not
measurably degrade learning relative to float64 on this task. The float64 loss
is genuinely lower, which is the approximation showing up where one would expect
it.

This is a four-feature binary logistic regression on 100 samples split five
ways. It says nothing about deep models.

\subsection{Consensus Cost and Behaviour Under Adversity}
\label{sec:eval-consensus}

\begin{table}[t]
\centering
\caption{Message cost against $n$, decomposed. Deliveries \emph{per height},
not broadcasts and not per run; 5 heights, 5\,ms one-way delay, no loss, $p=3$
proposers. The last column is $(n-1)(p+4n)$.}
\label{tab:scale}
\begin{tabular}{rrrrrrrrr}
\toprule
$n$ & $q$ & \textbf{Median} & \textbf{Prop.} & \textbf{Votes} &
\textbf{Polkas} & \textbf{Certs} & \textbf{Total} & \textbf{Closed form} \\
\midrule
5  & 4  & 0.26\,s  & 12  & 40   & 20   & 20   & 92    & 92 \\
10 & 7  & 0.46\,s  & 27  & 180  & 90   & 90   & 387   & 387 \\
20 & 14 & 2.04\,s  & 57  & 760  & 380  & 380  & 1577  & 1577 \\
50 & 34 & 26.45\,s & 147 & 4988 & 2450 & 2450 & 10094 & 9947 \\
\bottomrule
\end{tabular}
\end{table}

Table~\ref{tab:scale} shows message cost growing quadratically, not cubically:
deliveries per $n^2$ is flat, 3.68 to 4.04, as $n$ goes from 5 to 50. Our
earlier design re-broadcast an $O(n)$-sized vote set with every new vote, which
is $O(n^3)$ bytes per height.

We give the decomposition rather than a single total because a single total here
is indistinguishable from a formula, and would be right to distrust. It counts
\emph{deliveries}, so each of the $p$ proposals reaches $n-1$ peers; every node
sends a prevote and a precommit, gossips the polka it observed, and gossips the
certificate. That is $(n-1)(p + 4n)$, and at $n \le 20$ the measurement lands on
it exactly --- because at zero loss nothing is ever retried, so there is nothing
for the count to differ by. At $n = 50$ it does not: 10{,}094 against a predicted
9{,}947, the excess being sub-rounds that expired before a quorum arrived at a
26\,s median height. The count also scales linearly in heights, 144, 360 and 720
deliveries for 2, 5 and 10 heights at $n=5$, which is the check that it is
per-height rather than per-run.

The polka column is the price of the liveness fix in Section~\ref{sec:consensus}:
$n(n-1)$ deliveries per height, a quarter of the total. It buys termination out
of a locked/unlocked split that no number of sub-rounds resolves, and we report
it as its own line rather than folding it into the vote total, since a reader
comparing against a system without transferable polkas should see what it costs.

\begin{table}[t]
\centering
\caption{Behaviour under injected delay and loss, $n=10$, 25 heights across
three independent seeds per row. A height counts as committed only when
\emph{every} node holds it within a 15\,s deadline, so one node still syncing
marks the height failed; and because the driver re-attempts the same height, a
height that takes longer than the deadline is charged to every interval it
spans. These are the strictest readings of both words, and they matter for the
loss rows.}
\label{tab:adversity}
\begin{tabular}{rrcrl}
\toprule
\textbf{One-way delay} & \textbf{Loss} & \textbf{Heights committed} &
\textbf{Median} & \textbf{Per seed} \\
\midrule
5\,ms   & 0\%  & 75/75 & 0.35\,s  & 25/25, 25/25, 25/25 \\
50\,ms  & 0\%  & 75/75 & 0.50\,s  & 25/25, 25/25, 25/25 \\
150\,ms & 0\%  & 75/75 & 1.46\,s  & 25/25, 25/25, 25/25 \\
50\,ms  & 5\%  & 49/75 & 0.60\,s  & 3/25, 25/25, 21/25 \\
50\,ms  & 20\% & 3/75 & 15.01\,s & 1/25, 1/25, 1/25 \\
\bottomrule
\end{tabular}
\end{table}

Table~\ref{tab:adversity} is where we depart from the convention of reporting
only the configurations that worked. Delay is handled: 150\,ms
one-way costs latency and nothing else. Loss is a different matter, and reporting
it honestly required more care than we first gave it. Estimated from five
heights per configuration, the 5\%-loss figure moved between 5/5 and
3/5 across runs --- five samples is not a rate. Re-running at 25 heights across
three seeds gives 49/75, and the per-seed split of 3/25, 25/25, 21/25 shows the
variance is real rather than a sampling artefact: one seed in three collapses.

We initially attributed this to the sub-round budget $r_{\max}$, on the
reasoning that loss needs retries rather than longer windows, and a 10-height
sweep appeared to confirm it: 0.23 round-failure at $r_{\max} = 4$ against 0.00
at 8 and 16. Re-running that sweep at 25 heights per seed does not support the
conclusion. We measure failure rates of \textbf{0.48, 0.29 and 0.53} at
$r_{\max} = 4$, 8 and 16 --- better at 8 than at 4, but \emph{worse at 16 than
at either}, so the relationship is not monotone and, given the per-seed spread
above, the differences are not separable from noise with three seeds.

The non-monotonicity does have a mechanical explanation, and it is the reason we
report the budget as coupled rather than as a tuning knob. Attempt $r$ waits
$w 2^{r}$ (capped at $2^6 w$), so a larger budget also means later attempts are
individually longer; past the point where one window exceeds the height's
deadline, raising $r_{\max}$ buys unreachable sub-rounds rather than extra
attempts. Budget and window schedule have to be designed together, and we did
not. We therefore make \emph{no} tuning recommendation for lossy links beyond
the observation that none of the settings we tried makes 5\% loss reliable.

\paragraph{The deadline is the binding constraint, not the budget.}
That reasoning predicts something testable, so we tested it: hold $r_{\max}$
fixed at 8 and vary the per-height deadline instead. At 5\% loss, $n=10$ and
50\,ms one-way, a 15\,s deadline --- which the schedule above makes 5 to 6
attempts reachable within --- commits 22/24 heights, and a 45\,s deadline ---
6 to 8 reachable --- commits 24/24. Moving the deadline moves the outcome
while the budget is constant, which is the direct evidence that the budget was
never what bound. \texttt{experiments/e3\_deadline.py} prints the cumulative
window schedule beside the deadlines it tests, so the arithmetic above is
reproducible rather than asserted.

\paragraph{How much is any single loss figure worth?}
Because the simulator runs on real \texttt{asyncio} timers rather than a logical
clock, two runs with identical seeds do not see identical interleavings, and at
5\% loss the protocol sits close enough to a tipping point that the difference
shows. We therefore measured it rather than quoting one run: three repeats of
the same three seeds, 15 heights each. Commit rates across repeats were
0.82--1.00 (pooled 127/135), against 0.47--1.00 \emph{across seeds within} a
repeat. Seed-to-seed spread dominates run-to-run spread, which is the answer we
wanted --- the seed is doing more of the work than the scheduler --- but neither
is small, and a single number from a single run is a sample of a mixture. We
report pooled rates and ranges throughout this subsection for that reason, and
\texttt{make e3-repeat} regenerates the spread. It must be run on an idle
machine: an earlier attempt, taken while another experiment was running, gave
0.71/0.40/0.33 in the order the other job got heavier, and most of that
``run-to-run spread'' was the host rather than the protocol.

At 20\% loss no budget helps at all, and this is a limit of the design rather
than a parameter to tune. A lossy link needs a retransmitting transport
underneath; this protocol does not compensate for one.

\subsection{Adversarial Scenarios}
\label{sec:eval-adversarial}

Three live adversaries at $n=10$, $f=3$, $q=7$, three heights each. An
\textbf{equivocating voter} and an \textbf{equivocating proposer} are each
detected by all nine honest nodes, with converged blacklists, \emph{zero} false
positives, and all three heights still committing. A \textbf{flooder} is
rate-limited on 7{,}895 messages without stalling the network; it is not
blacklisted, since flooding is not equivocation and there is nothing to prove
against it.

\begin{table}[t]
\centering
\caption{Poisoning by nodes that train correctly on committed-but-adversarial
data, 40 heights, $n=10$, trim $=250$ per mille. Median over three data seeds,
with the range across seeds; clean-run accuracy 0.947 (0.928--0.958).
\textbf{Mag.} is the feature scaling the adversary committed to, chosen by
sweeping and keeping the most damaging setting that stays inside the clipping
bound --- so no poisoned step is ever discarded as unprovable. The last column
is the trimmed mean's own robustness precondition, that at least $f$ values are
trimmed per side; it is a property of the rule's configuration, not of the
clipping bound, which holds in every row by construction.}
\label{tab:poison}
\begin{tabular}{clcccrc}
\toprule
\textbf{Poisoned} & \textbf{Rule} & \textbf{Mag.} & \textbf{Accuracy} &
\textbf{Lost} & $\|\mu - \mu_{\rm clean}\|_\infty$ & \textbf{Trim $\ge f$} \\
\midrule
\multirow{3}{*}{2 of 10}
 & FedAvg       & 1.0  & 0.842 (0.761--0.914) & $-0.117$ & 857{,}677 & --- \\
 & Trimmed mean & 0.5  & 0.931 (0.928--0.933) & $-0.017$ & 342{,}623 & yes \\
 & Median       & 0.25 & 0.933 (0.925--0.936) & $-0.014$ & 316{,}982 & yes \\
\midrule
\multirow{3}{*}{3 of 10}
 & FedAvg       & 1.0  & 0.831 (0.700--0.892) & $-0.128$ & 1{,}001{,}866 & --- \\
 & Trimmed mean & 1.0  & 0.844 (0.739--0.881) & $-0.114$ & 791{,}064 & \textbf{no} \\
 & Median       & 0.25 & 0.906 (0.903--0.914) & $-0.033$ & 568{,}312 & yes \\
\bottomrule
\end{tabular}
\end{table}

Table~\ref{tab:poison} reports the attack the circuit cannot stop. Each poisoned
node commits to a dataset of large-magnitude features with the labels a correct
model would not assign, registers its Merkle root, and then trains \emph{honestly}
on it. Every resulting proposal carries a valid proof, because it is a correct
gradient step on committed data. Only the aggregation rule can help.

\paragraph{The adversary is calibrated, not merely loud.}
An earlier version of this experiment let the poisoned nodes propose whatever
their data produced and discarded the steps that exceeded $\delta_{\max}$ as
unprovable, counting them. That is honest but it is not the strongest attack:
a real adversary does not throw away its own rounds. It commits, once and in
advance, to the dataset whose steps do the most damage \emph{while remaining
provable}. We model that by sweeping the feature scaling and keeping, for each
rule, the most damaging setting that never trips the bound --- the \textbf{Mag.}
column. No poisoned step is discarded in any row of Table~\ref{tab:poison}.

The calibration is not a formality. Against FedAvg the adversary picks the
largest scaling it can, because a mean has no defence against magnitude. Against
trimmed mean and median it deliberately picks a \emph{smaller} one (0.5 and
0.25): an extremal contribution is exactly what an order statistic discards, so
the better attack is the one that stays inside the kept range and pulls it. This
roughly triples the damage the earlier attack did to trimmed mean at 2 poisoned
nodes, from 0.6\% to 1.7\%.

\medskip
Two results survive that strengthening. First, at 2 poisoned nodes the trimming
threshold covers the adversaries and trimmed mean holds: FedAvg loses 11.7\%
accuracy where trimmed mean loses 1.7\%. Second, at 3 poisoned nodes,
$\lfloor 10 \cdot 250/1000 \rfloor = 2$ values are trimmed per side while
$f = 3$, so the robustness precondition is violated --- and trimmed mean loses
11.4\%, essentially matching plain FedAvg. It fails exactly where its own bound
says it will. Our implementation emits a runtime warning on that configuration
and no other, and the measurement confirms the warning is load-bearing rather
than decorative. Median, which has no trimming threshold, degrades gracefully in
both.

The seed ranges deserve reading alongside the medians. FedAvg's accuracy varies
from 0.761 to 0.914 across three seeds at 2 poisoned nodes, which is wider than
the gap it is being compared against; the robust rules vary by under a point.
That the \emph{unprotected} rule is also the least predictable one is
unsurprising, but it means the FedAvg rows should be read as ``much worse, with
high variance'' rather than as a specific number.

\subsection{Offline Auditability}
\label{sec:eval-audit}

We ran 7 nodes for 3 heights in strict mode, published node 0's audit log as
plain JSON (41.6\,kB, 13.9\,kB per height), and
re-verified it from a cold start with a fresh proof backend sharing no state
with the run: all certificates re-checked against the member set, all proofs
re-checked against the model as of their height and the proposer's registered
root, all aggregates recomputed, all chain links recomputed. Verification took
0.37\,s, or 0.12\,s per height.

\begin{table}[t]
\centering
\caption{Tampering with the published log. Nine attacks, nine rejections.}
\label{tab:tamper}
\begin{tabular}{ll}
\toprule
\textbf{Tampering} & \textbf{Rejected because} \\
\midrule
Model edited                        & model does not match its hash \\
Model edited, hash recomputed       & recomputed aggregate does not match \\
Certificate one vote short          & certificate lacks a quorum \\
Certificate padded with a duplicate voter & certificate lacks a quorum \\
Precommits relabelled as prevotes   & certificate lacks a quorum \\
A proposal dropped from the set     & recomputed aggregate does not match \\
Proofs swapped between proposers    & bad proposal signature \\
Proofs swapped \emph{and re-signed} & proof does not verify \\
A height spliced out                & broken parent link \\
\bottomrule
\end{tabular}
\end{table}

Table~\ref{tab:tamper} gives the negative results, which are the point of the
experiment: a checker that accepts everything checks nothing. The last two rows
are a pair. Swapping two proposals' proofs is caught by the signature check,
because the signature covers the proof --- which means the pairing check is never
exercised. Re-signing the swapped proposals with the genuine proposers' keys
makes every message well-formed and leaves proof verification as the only
remaining defence; it rejects.

The first row is one this experiment added. Writing the tampering suite is what
found the corresponding defect: our offline verifier recomputed the aggregate
and compared \emph{its} hash, but never checked the entry's human-readable model
field against that hash. An edited model therefore passed a full audit, because
the field a reader actually consumes was not the field being checked. The suite
reported \textsc{accepted} where it expected a rejection, which is the only
reason we found it. We record this because it is the argument for negative
tests: every other row passed from the first run, and a suite containing only
those rows would have certified a verifier with a hole in it.

% ─────────────────────────────────────────────────────────────────────────────
\section{Discussion and Limitations}
\label{sec:limitations}

\paragraph{The proof does not make the update private.}
This is the most important limitation and it is structural. The zk-SNARK hides
the training data. It does not hide $\mathbf{w}_{\rm new}$, which is public by
construction: every node needs it to aggregate. Gradient-inversion
attacks~\cite{zhu2019deepleakage,geiping2020inverting} reconstruct training
samples from exactly this quantity. \emph{Publishing a per-node update weakens
privacy regardless of the proof.} Composing this design with differential
privacy or secure aggregation is future work; claiming privacy without it would
be wrong.

\paragraph{Is the zero-knowledge property doing any work?}
It is a fair challenge, and the honest answer is: less than the name suggests,
but not nothing. We publish $\mathbf{w}_{\rm new}$, and we cite the
gradient-inversion literature to say that publishing it leaks information about
the batch. So the protocol's \emph{integrity} requirement is a succinct
argument, not a zero-knowledge one, and a reader is right to ask why not use a
transparent, setup-free system --- which would additionally delete our
trusted-setup limitation.

Two things the hiding property still buys. First, the witness is the batch
\emph{and} its Merkle path: what stays hidden is which $k$ of the node's
committed samples were used at this height, and thereby the rest of the
committed dataset, which no update reveals. A verifier learns that some
committed batch was used, not which. Second, and more prosaically, Groth16 gives
a constant-size proof and millisecond verification, and in a gossip protocol
every node verifies every proposal at every height --- $O(n)$ verifications per
node per height. A transparent system with proofs in the tens of kilobytes and
verification in the tens of milliseconds changes the message-cost table in
Section~\ref{sec:eval-consensus} substantially, and we did not want to claim a
consensus cost we had measured with a cheaper primitive.

That is a defensible engineering choice rather than a cryptographic necessity,
and we state it as one. A transparent proof system is the natural next step for
a deployment that would rather lose proof size than keep a ceremony, and nothing
in the protocol above depends on Groth16 specifically.

\paragraph{The circuit cannot distinguish honest data from poisoned data.}
It proves that the update \emph{is} the gradient step on the \emph{committed}
data. A node that commits to a poisoned dataset and trains correctly produces a
valid proof, as Section~\ref{sec:eval-adversarial} demonstrates. No circuit of
this form can do otherwise; defending it requires data-level attestation.

\paragraph{Nothing binds a batch to a height.}
The circuit proves that the update is a correct step on \emph{some} batch whose
digest is in the committed tree. It does not constrain \emph{which}. A node is
therefore free to train on its most damaging committed batch at every height,
and an adversary certainly would. Honest nodes in our implementation walk their
batches round-robin, and the experiment of
Section~\ref{sec:eval-adversarial} does the same, so the poisoning numbers
\emph{understate} this adversary: it commits one poisoned shard and then replays
its worst batch indefinitely.

The fix is small and we describe it rather than claim it. Expose the Merkle leaf
index as a public signal, constrain it to the path bits already in the circuit,
and have the verifier --- which already supplies every other public input from
its own state --- require it to equal a deterministic schedule, for instance
$H(h \,\|\, \text{proposer}) \bmod 2^d$. That makes batch selection a function of
the height rather than a choice, at the cost of one public input and a handful
of constraints. We did not implement it, and the omission is a gap in coverage
rather than in soundness: the proof still establishes exactly what it claims.

\paragraph{Robustness is conditional, and the condition is checkable.}
Coordinate-wise trimmed mean bounds adversarial influence only when at least $f$
values are trimmed per side. With few proposals per height the trim count rounds
down and the bound becomes vacuous. Our implementation says so at runtime rather
than staying silent, and Table~\ref{tab:poison} measures the consequence.

\paragraph{Trusted setup.}
Per-circuit, as Groth16 requires. Sound if one phase-2 contributor is honest;
forgeable by the operator in the offline build mode, which is labelled as such.

\paragraph{Security level.}
bn128 offers roughly 100 bits of security after the extended tower number field
sieve~\cite{kim2016extnfs,barbulescu2019updating}, not the 128 sometimes quoted.

\paragraph{Scale.}
Four features and a batch of eight. Constraint count grows roughly linearly in
$n \cdot k$; extrapolating from 19{,}402 constraints to a convolutional network
is not supported by anything measured here, and the proof-of-training
literature~\cite{garg2023zkpot,abbaszadeh2024kaizen} is the appropriate
reference for that regime.

\paragraph{Approximation.}
The five-piece sigmoid has maximum error $0.0117$. It is adequate for logistic
regression and would need revisiting --- most plausibly via lookup arguments --- for
anything else.

\paragraph{Availability.}
A partition holding more than $f$ nodes stops progress. This is the price of
safety under partial synchrony~\cite{fischer1985flp} rather than a defect, but a
deployment must expect it. Sustained packet loss above a few percent degrades
progress sharply (Table~\ref{tab:adversity}).

\paragraph{Transport.}
Our WebSocket transport binds node identities to addresses and drops unknown
senders, but the channel is plain TCP. Every message carries its own signature,
so a network adversary can delay and drop but not forge; given the privacy
limitation above, however, a deployment should terminate TLS.

% ─────────────────────────────────────────────────────────────────────────────
\section{Conclusion}

We have presented a decentralized federated learning system in which each model
update carries a zk-SNARK binding it to a committed dataset and to the verifier's
own view of the global model, and in which nodes agree on the aggregate of the
verified updates through a leaderless two-phase BFT protocol with a stated fault
model and a proved agreement property. Misbehaviour produces transferable
evidence; the committed history is re-verifiable offline by a stranger, and we
show nine distinct tamperings being rejected.

The recurring lesson of this work is that each of these components fails quietly
when built approximately. A circuit that omits range checks and remainder
constraints proves nothing while appearing to prove something. A single-phase
consensus deadlocks on any network with real delay while passing every test on
loopback. A two-phase one deadlocks too, on a split between locked and unlocked
correct nodes that no number of sub-rounds resolves, unless the quorum that
licenses a lock can itself travel. A bound that holds for the honest witness
because a comparator chain determines it is not a constraint. Detecting an
equivocating proposer without retracting what it already had accepted splits the
network permanently. An offline auditor that checks a model's hash but not the
model field accepts a forgery. Almost all of these were present in our own
earlier design, and every one was found by an experiment or a proof obligation
built to look for it rather than by inspection --- several of them only when we
tried to write down why the code was correct and found that we could not.

Future work: composing with differential privacy or secure aggregation to
address the update-privacy limitation, which we regard as the most consequential
gap; lookup-argument activations to move beyond logistic regression; adaptive
quorum under churn; and a transport that survives the loss rates at which the
present protocol does not.

% ─────────────────────────────────────────────────────────────────────────────
% Single bibliography.  Entries live in paper/references.bib; compile with
% bibtex.  (The submitted draft carried two \thebibliography blocks whose
% second contained placeholder arXiv identifiers -- that duplication is gone.)
\bibliographystyle{splncs04}
\bibliography{references}

\end{document}
